
Reliable uncertainty quantification (UQ) for large language models (LLMs) is of practical importance for systems that must detect hallucinations, route queries to human reviewers, or abstain from low-confidence predictions. Token-probability-based approaches—using the model's own predicted logits as a confidence proxy—are computationally simple and require only a single forward pass. Sampling-based semantic methods were proposed as improvements that operate over semantic content rather than surface-level token distributions.

Semantic entropy (SE), introduced by Farquhar et al. (2024) and published in \textit{Nature}, clusters N stochastic responses into semantic equivalence classes via bidirectional NLI and computes entropy over cluster probability mass. Published evaluations report SE achieving AUROC of approximately 0.72–0.79 on TriviaQA, NaturalQuestions, BioASQ, and SQuAD, consistently exceeding token-probability (approximately 0.67 AUROC on TriviaQA). This superiority has been taken as established and has motivated several subsequent methods: Semantic Entropy Probes (SEPs; Kossen et al., 2024), which approximate SE from a single forward pass; Kernel Language Entropy (KLE; Nikitin et al., 2024), which generalizes SE via von Neumann entropy; and SelfCheckGPT (Manakul et al., 2023), which measures cross-sample NLI consistency. However, none of these evaluations explicitly reports the sampling diversity—specifically, the proportion of queries for which all N samples are semantically identical—as a diagnostic for the validity of clustering-based methods.

The present study was designed to test the Epistemic Geometry Scaling Hypothesis (EGSH), which predicted that the AUROC advantage of SE and KLE over token-probability would increase with model scale from 8B to 70B parameters in the Llama-3 base model family. As a prerequisite, the existence of an SE advantage at 8B was verified first. This existence check failed: SE achieves an AUROC of 0.4735 on TriviaQA and 0.5524 on NaturalQuestions when applied to Llama-3-8B-Base under the standard N=10 sampling protocol at temperature=1.0—both values substantially below token-probability on the same datasets.

The experiment records a degenerate_fraction of 0.894 on TriviaQA and 0.848 on NaturalQuestions: the proportion of queries for which all 10 stochastic samples from Llama-3-8B-Base are assigned to a single semantic cluster by the DeBERTa-large-mnli NLI model. When K=1 for a query, SE entropy is exactly zero regardless of whether the query is answered correctly or incorrectly. With 89\% of queries at K=1, SE scores cluster near zero and carry no discriminative signal. This is not an implementation error: the SE mechanism activates correctly (mean_k=9.884 < N=10, confirming that clustering runs), but the model's output distribution lacks the diversity that clustering-based entropy estimation requires to function as a UQ signal.

The paper is organized as follows. Section 2 reviews related work. Section 3 describes the experimental methodology. Section 4 presents results. Section 5 discusses implications and limitations. Section 6 concludes.

